% 原书 p.10–p.12（PDF p.18–p.20）
% >>> p018
\section{$J$-全纯曲线（第一部分）（$J$-Holomorphic Curve, Part One）}

\begin{keytakeaways}
\begin{itemize}
  \item 几何图像：$J$-全纯曲线是 Riemann 曲面到近复流形的映射，其微分与复结构交换；
        它正是复几何中全纯曲线在近复流形中的推广。
  \item 数学工具：Cauchy--Riemann 算子 $\bar\partial_J$、能量的积分表达式
        $E(u)=\int_\Sigma u^*\omega$、以及由 $g_J$ 给出的面积泛函 (1.10)。
  \item 本节结论：$J$-全纯曲线恰是能量泛函的下界等号情形，即 $E(u)\le\Area(u)$ 恒成立，
        且等号成立当且仅当 $u$ 是 $J$-全纯的（命题 1.4.1）。
  \item 推论：$J$-全纯曲线是极小曲面（注 1.4.1），且其能量只依赖同调类
        $u_*[\Sigma]\in H_2(M;\Z)$ 与 $[\omega]$ 的配对，故有拓扑下界。
\end{itemize}
\end{keytakeaways}

与复几何中最重要的课题之一是研究复流形中的曲线类似，人们在近复流形中定义类似的曲线。
正是 Gromov 的天才工作 [99] 在 1980 年代把这一套机器带进了辛几何，使这一数学分支发生了
革命性的变化，其巨大影响甚至延续至今。

\begin{wdefn}{1.4.1}
从 Riemann 曲面 $(\Sigma_g,j)$（可以带边界）到近复流形 $(M,J)$ 的光滑映射 $\varphi$
称为 \textbf{$J$-全纯的}（$J$-holomorphic），若它满足
\[
  \varphi_*\circ j=J\circ\varphi_* .
\]
\end{wdefn}

为消化定义 1.4.1，让我们把源与靶都取为 $(\C,\sqrt{-1})$，其中 $\sqrt{-1}$ 只是指那个（近）复结构
——即乘以 $\sqrt{-1}$。我们断言定义 1.4.1 中的条件等价于经典的 Cauchy--Riemann 方程。
下面的例子对此作了阐述。

\begin{wexam}{1.4.1}
光滑映射 $\varphi:(\C,\sqrt{-1})\to(\C,\sqrt{-1})$ 按定义是 $\sqrt{-1}$-全纯的，
若 $\sqrt{-1}\circ\varphi_*=\varphi_*\circ\sqrt{-1}$。让我们借助 $z=x+\sqrt{-1}\,y$ 把 $z\in\C$
与 $(x,y)\in\R^2$ 等同起来，并写 $\varphi(x,y)=u(x,y)+\sqrt{-1}\,v(x,y)$，则在这些坐标下，
我们有如下关系：
\[
  \varphi_*=\begin{pmatrix}
    \dfrac{\partial u}{\partial x} & \dfrac{\partial v}{\partial x}\\[2mm]
    \dfrac{\partial u}{\partial y} & \dfrac{\partial v}{\partial y}
  \end{pmatrix},
  \qquad
  \sqrt{-1}=\begin{pmatrix}0&1\\-1&0\end{pmatrix},
\]
以及
\[
  \begin{aligned}
    \sqrt{-1}\circ\varphi_*
    &=\begin{pmatrix}0&1\\-1&0\end{pmatrix}
      \cdot\begin{pmatrix}
        \dfrac{\partial u}{\partial x} & \dfrac{\partial v}{\partial x}\\[2mm]
        \dfrac{\partial u}{\partial y} & \dfrac{\partial v}{\partial y}
      \end{pmatrix}
     =\begin{pmatrix}
        \dfrac{\partial u}{\partial y} & \dfrac{\partial v}{\partial y}\\[2mm]
        -\dfrac{\partial u}{\partial x} & -\dfrac{\partial v}{\partial x}
      \end{pmatrix},\\
    \varphi_*\circ\sqrt{-1}
    &=\begin{pmatrix}
        \dfrac{\partial u}{\partial x} & \dfrac{\partial v}{\partial x}\\[2mm]
        \dfrac{\partial u}{\partial y} & \dfrac{\partial v}{\partial y}
      \end{pmatrix}
      \cdot\begin{pmatrix}0&1\\-1&0\end{pmatrix}
     =\begin{pmatrix}
        -\dfrac{\partial v}{\partial x} & \dfrac{\partial u}{\partial x}\\[2mm]
        -\dfrac{\partial v}{\partial y} & \dfrac{\partial u}{\partial y}
      \end{pmatrix}.
  \end{aligned}
\]
于是 $\sqrt{-1}$-全纯条件等价于
$\dfrac{\partial u}{\partial x}=\dfrac{\partial v}{\partial y}$ 与
$-\dfrac{\partial u}{\partial y}=\dfrac{\partial v}{\partial x}$。\hfill$//$
\end{wexam}
% [Tutor Note: 这里 $\varphi_*$ 是以 $\{x,y\}$ 与 $\{u,v\}$ 为基的 Jacobi 矩阵，故第 1 列为
% $\partial\varphi/\partial x=(\partial_xu,\partial_xv)$、第 2 列为
% $\partial\varphi/\partial y=(\partial_yu,\partial_yv)$。易错点：$\sqrt{-1}$ 在这两个切空间上
% 都是同一个 $2\times2$ 矩阵，但两边的定向约定一旦相反（例如把靶上的 $J$ 取为
% $(0,-1;1,0)$），得到的方程会变成反全纯方程 $u_x=-v_y$、$u_y=v_x$。另外，
% 「$\varphi_*\circ j=J\circ\varphi_*$」是复几何中 $\bar\partial_J\varphi=0$ 的坐标无关写法；
% 这里 $j$、$J$ 分别记源与靶的近复结构，已在定义 1.4.1 中分清。]

% >>> p019
一旦给 $M$ 装备了辛结构（回忆由命题 1.3.1，总存在 $\omega$-相容的近复结构 $J$），我们就能对
这些 $J$-全纯曲线做定量研究。为强调辛结构，让我们把所有数据装进一个三元组 $(M,\omega,J)$。
此外，为了与其他文献保持一致，$(M,\omega,J)$ 中的 $J$-全纯曲线通常记作
$u:(\Sigma,j)\to(M,\omega,J)$。这里 $\Sigma=\Sigma_g$，而在全书中我们只关心亏格 $g=0$ 的曲线。

借用微分几何中的概念，对 $J$-全纯曲线 $u:(\Sigma,j)\to(M,\omega,J)$，人们定义如下两个量来
度量它的大小。
\begin{enumerate}[label=(\roman*)]
  \item $u$ 的\textbf{能量}（energy）由下式给出
    \begin{equation}
      E(u):=\int_\Sigma u^*\omega .
      \tag{1.9}
    \end{equation}
  \item $u$ 相对于 $g_J$ 的\textbf{面积} $\Area(u)$ 局部地由下式定义
    \begin{equation}
      \sqrt{g_J(X,X)g_J(Y,Y)-g_J(X,Y)^2}
      \tag{1.10}
    \end{equation}
    其中 $\{X,Y\}$ 是 $u_*T\Sigma$ 的任意一组基。特别地，选取 $T\Sigma$ 的一组基
    $\{e_1,e_2\}$ 使得 $g_J(u_*e_1,u_*e_2)=0$，则这个局部面积为
    $\sqrt{g_J(u_*e_1,u_*e_1)\,g_J(u_*e_2,u_*e_2)}$。
\end{enumerate}

下面这个结果比较了这两个量。

\begin{wprop}{1.4.1}
对任意曲线 $u:(\Sigma,j)\to(M,\omega,J)$，都有 $E(u)\le\Area(u)$。当 $u$ 是 $J$-全纯曲线时
等号成立。
\end{wprop}

\begin{wproof}
在 $T\Sigma$ 上取一组基 $\{e_1,e_2\}$，使得 $g_J(u_*e_1,u_*e_2)=0$。则
\[
  (u^*\omega)(e_1,e_2)=\omega(u_*e_1,u_*e_2)=\omega(Ju_*e_1,Ju_*e_2)=g_J(Ju_*e_1,u_*e_2).
\]
由 Cauchy--Schwarz 不等式，
\[
  \begin{aligned}
    g_J(Ju_*e_1,u_*e_2)
    &\le\sqrt{g_J(Ju_*e_1,Ju_*e_1)\,g_J(u_*e_2,u_*e_2)}\\
    &=\sqrt{\omega(Ju_*e_1,-u_*e_1)\,g_J(u_*e_2,u_*e_2)}\\
    &=\sqrt{g_J(u_*e_1,u_*e_1)\,g_J(u_*e_2,u_*e_2)} .
  \end{aligned}
\]
回忆 Cauchy--Schwarz 不等式取等当且仅当 $Ju_*e_1=u_*e_2$。现在，若 $u$ 是 $J$-全纯的，
则调整 $\Sigma$ 上的基 $\{e_1,e_2\}$ 使得 $e_2=je_1$，于是
$Ju_*e_1=u_*je_1=u_*e_2$，这有助于在 Cauchy--Schwarz 不等式中实现等号。需要强调的是，
这一调整之所以可行，是因为我们可以从任意满足 $e_2=je_1$ 的基 $\{e_1,e_2\}$ 出发，而它自动
满足本证明开头所提的要求，即
$g_J(u_*e_1,u_*e_2)=g_J(u_*e_1,u_*je_1)=g_J(u_*e_1,Ju_*e_1)=\omega(u_*e_1,-u_*e_1)=0$，
其中第二个等号来自「$u$ 是 $J$-全纯的」这一假设。
\end{wproof}
% [Tutor Note: 此处 $g_J(v,w)=\omega(v,Jw)$ 是 (1.8) 的约定，于是
% $g_J(Ju_*e_1,u_*e_2)=\omega(Ju_*e_1,Ju_*e_2)$ —— 注意 $J$ 只作用在第一个槽上，
% 第二个槽里没有 $J$；而 $\omega(Ju_*e_1,Ju_*e_2)=\omega(u_*e_1,u_*e_2)$ 依据的正是 (1.7)(i)。
% 再注意 $g_J(Ju_*e_1,Ju_*e_1)=\omega(Ju_*e_1,J(Ju_*e_1))=\omega(Ju_*e_1,-u_*e_1)$ 用了
% $J^2=-\mathbf{1}$；最后 $\omega(Ju_*e_1,-u_*e_1)=\omega(u_*e_1,Ju_*e_1)=g_J(u_*e_1,u_*e_1)$
% 又用了一次 (1.8)。这几步符号稍有疏忽就会得到 $-E$ 而非 $E$，是本节最易出错之处。]

\begin{wrem}{1.4.1}
命题 1.4.1 是局部地验证的，因此一条 $J$-全纯曲线必然是一个极小曲面，即在局部极小化
(1.10) 所定义的那个面积泛函。
\end{wrem}

% >>> p020
\begin{wexer}{1.4.1}（[229] 中的注 2.5）
若 $J$ 只是 $\omega$-驯顺的（见定义 1.3.2），我们同样可以在 $M$ 上定义 Riemann 度量
\[
  g_J(v,w):=\frac{\omega(v,Jw)+\omega(w,Jv)}{2}.
\]
在此情形下，请检验 $J$-全纯曲线 $u$ 是否仍为极小曲面。换言之，请检验
\[
  E(u)\le\Area(u),
\]
其中 $E(u)$ 与 $\Area(u)$ 分别由 (1.9) 与 (1.10) 定义。
\end{wexer}
% [Tutor Note: 驯顺情形下 $g_J$ 需要对称化，因为 $J$ 与 $\omega$ 不相容时 $\omega(v,Jw)$ 本身
% 未必对称；对称化后 $g_J$ 仍正定（由 (1.7)(ii)），但一般不再 $J$-不变，即
% $g_J(Jv,Jw)\neq g_J(v,w)$，这会使命题 1.4.1 的 Cauchy–Schwarz 估计多出一项误差，
% 因此本题的答案与相容情形不同：等号不再必然成立，需利用 $g_J$ 与
% $\omega(\cdot,J\cdot)$ 的偏差来估计。]
